% GENERATED FILE. Edit canonical lessons, then run npm run generate:books.
% canonical-source-hash: fnv1a64:2d70819b0aad6479
% canonical-lessons: LA-C409-edictum, LA-C409-decretum, LA-C409-praeceptum, LA-C409-legitimus, LA-C409-licitus

\chapter{Edicts and Rules}
\label{ch:la-edicts-and-rules}
\hlchaptermodality{409}

\begin{chapteropening}
\textbf{By the end of this chapter:} \emph{I can talk about an edict, a decree and a rule, and say that something is lawful or allowed.}

Complete the last lesson of chapter 409: I can talk about an edict, a decree and a rule, and say that something is lawful or allowed.
\end{chapteropening}

\section[ēdictum, ēdictī]{ēdictum, ēdictī — an edict, a proclamation}

\label{lesson:LA-C409-edictum}



\begin{quote}
Before the new one: say the Latin for a lash, then the Latin for to be beaten.
\end{quote}

\subsection*{You'll want to know: ēdictum, ēdictī}
\textbf{ēdictum, ēdictī} — \textquotedblleft{}an edict, a proclamation\textquotedblright{}.

\begin{quote}
Praetor ēdictum prōposuit.
\end{quote}

\begin{quote}
\emph{The praetor published his edict.}
\end{quote}

\emph{ēdictum, ēdictī} n., from \emph{ēdīcō}, to speak out. Each praetor posted an edict saying how he would run the courts. English \emph{edict}.

\begin{cousinweb}
\emph{ēdictum} is the participle of \emph{ēdīcō} \textquotedblleft{}to speak out, to proclaim\textquotedblright{} (EtymDB; Lewis \& Short), from \emph{ē} and \emph{dīcō} \textquotedblleft{}to say\textquotedblright{}. English \emph{edict} comes from it.
\end{cousinweb}

\begin{grammarlens}[title={What you've built}]
One more word for telling, explaining and describing at B1.
\end{grammarlens}

\subsection*{Your turn}
\begin{itemize}
\raggedright
  \item \emph{Say it:} \emph{ēdictum, ēdictī}
  \item \emph{Say it:} \emph{ēdictum, ēdictī}, once more
  \item \emph{Say it:} a sentence of your own with \emph{ēdictum, ēdictī}
  \item \emph{Recall:} say the Latin for a lash, then the Latin for to be beaten, then say \emph{ēdictum, ēdictī} again
\end{itemize}

\begin{tcolorbox}[breakable,colback=teal!4,colframe=teal!35!black,title={Before you move on}]
What does ēdictum, ēdictī mean? (\textquotedblleft{}An edict, a proclamation\textquotedblright{}.) Say it once more.
\end{tcolorbox}
\section[dēcrētum, dēcrētī]{dēcrētum, dēcrētī — a decree, a decision}

\label{lesson:LA-C409-decretum}



\begin{quote}
Before the new one: say the Latin for to be beaten, then the Latin for an edict.
\end{quote}

\subsection*{You'll want to know: dēcrētum, dēcrētī}
\textbf{dēcrētum, dēcrētī} — \textquotedblleft{}a decree, a decision\textquotedblright{}.

\begin{quote}
Senātūs dēcrētum factum est.
\end{quote}

\begin{quote}
\emph{A decree of the senate was passed.}
\end{quote}

\emph{dēcrētum, dēcrētī} n., the participle of \emph{dēcernō}, to decide. English \emph{decree}.

\begin{cousinweb}
\emph{dēcrētum} is the participle of \emph{dēcernō} \textquotedblleft{}to decide, to decree\textquotedblright{} (EtymDB; Lewis \& Short), from \emph{dē} and \emph{cernō} \textquotedblleft{}to sift\textquotedblright{}. English \emph{decree} comes from it through French.
\end{cousinweb}

\begin{grammarlens}[title={What you've built}]
One more word for telling, explaining and describing at B1.
\end{grammarlens}

\subsection*{Your turn}
\begin{itemize}
\raggedright
  \item \emph{Say it:} \emph{dēcrētum, dēcrētī}
  \item \emph{Say it:} \emph{dēcrētum, dēcrētī}, once more
  \item \emph{Say it:} a sentence of your own with \emph{dēcrētum, dēcrētī}
  \item \emph{Recall:} say the Latin for to be beaten, then the Latin for an edict, then say \emph{dēcrētum, dēcrētī} again
\end{itemize}

\begin{tcolorbox}[breakable,colback=teal!4,colframe=teal!35!black,title={Before you move on}]
What does dēcrētum, dēcrētī mean? (\textquotedblleft{}A decree, a decision\textquotedblright{}.) Say it once more.
\end{tcolorbox}
\section[praeceptum, praeceptī]{praeceptum, praeceptī — a rule, a precept, an instruction}

\label{lesson:LA-C409-praeceptum}



\begin{quote}
Before the new one: say the Latin for an edict, then the Latin for a decree.
\end{quote}

\subsection*{You'll want to know: praeceptum, praeceptī}
\textbf{praeceptum, praeceptī} — \textquotedblleft{}a rule, a precept, an instruction\textquotedblright{}.

\begin{quote}
Praecepta medicī servāvit.
\end{quote}

\begin{quote}
\emph{He followed the doctor's instructions.}
\end{quote}

\emph{praeceptum, praeceptī} n., the participle of \emph{praecipiō}. English \emph{precept}.

\begin{cousinweb}
\emph{praeceptum} is the participle of \emph{praecipiō} \textquotedblleft{}to instruct\textquotedblright{} (EtymDB; Lewis \& Short), from \emph{prae} and \emph{capiō}. English \emph{precept} comes from it.
\end{cousinweb}

\begin{grammarlens}[title={What you've built}]
One more word for telling, explaining and describing at B1.
\end{grammarlens}

\subsection*{Your turn}
\begin{itemize}
\raggedright
  \item \emph{Say it:} \emph{praeceptum, praeceptī}
  \item \emph{Say it:} \emph{praeceptum, praeceptī}, once more
  \item \emph{Say it:} a sentence of your own with \emph{praeceptum, praeceptī}
  \item \emph{Recall:} say the Latin for an edict, then the Latin for a decree, then say \emph{praeceptum, praeceptī} again
\end{itemize}

\begin{tcolorbox}[breakable,colback=teal!4,colframe=teal!35!black,title={Before you move on}]
What does praeceptum, praeceptī mean? (\textquotedblleft{}A rule, a precept, an instruction\textquotedblright{}.) Say it once more.
\end{tcolorbox}
\section[lēgitimus, lēgitima, lēgitimum]{lēgitimus, lēgitima, lēgitimum — lawful, legitimate; proper}

\label{lesson:LA-C409-legitimus}



\begin{quote}
Before the new one: say the Latin for a decree, then the Latin for a rule.
\end{quote}

\subsection*{You'll want to know: lēgitimus, lēgitima, lēgitimum}
\textbf{lēgitimus, lēgitima, lēgitimum} — \textquotedblleft{}lawful, legitimate; proper\textquotedblright{}.

\begin{quote}
Hoc nōn est lēgitimum.
\end{quote}

\begin{quote}
\emph{This is not lawful.}
\end{quote}

\emph{lēgitimus, lēgitima, lēgitimum}, built on \emph{lēx, lēgis}. English \emph{legitimate}.

\begin{cousinweb}
\emph{lēgitimus} is built on \emph{lēx, lēgis} \textquotedblleft{}a law\textquotedblright{} (EtymDB; Lewis \& Short). English \emph{legitimate} comes from it.
\end{cousinweb}

\begin{grammarlens}[title={What you've built}]
One more word for telling, explaining and describing at B1.
\end{grammarlens}

\subsection*{Your turn}
\begin{itemize}
\raggedright
  \item \emph{Say it:} \emph{lēgitimus, lēgitima, lēgitimum}
  \item \emph{Say it:} \emph{lēgitimus, lēgitima, lēgitimum}, once more
  \item \emph{Say it:} a sentence of your own with \emph{lēgitimus, lēgitima, lēgitimum}
  \item \emph{Recall:} say the Latin for a decree, then the Latin for a rule, then say \emph{lēgitimus, lēgitima, lēgitimum} again
\end{itemize}

\begin{tcolorbox}[breakable,colback=teal!4,colframe=teal!35!black,title={Before you move on}]
What does lēgitimus, lēgitima, lēgitimum mean? (\textquotedblleft{}Lawful, legitimate; proper\textquotedblright{}.) Say it once more.
\end{tcolorbox}
\section[licitus, licita, licitum]{licitus, licita, licitum — allowed, lawful}

\label{lesson:LA-C409-licitus}



\begin{quote}
Before the new one: say the Latin for a rule, then the Latin for lawful.
\end{quote}

\subsection*{You'll want to know: licitus, licita, licitum}
\textbf{licitus, licita, licitum} — \textquotedblleft{}allowed, lawful\textquotedblright{}.

\begin{quote}
Quod nōn licitum est, nōn faciō.
\end{quote}

\begin{quote}
\emph{What is not allowed, I do not do.}
\end{quote}

\emph{licitus, licita, licitum}, the participle of \emph{licet}, it is allowed. Its opposite is \emph{illicitus}. English \emph{licit} and \emph{illicit}.

\begin{cousinweb}
\emph{licitus} is the participle of \emph{licet} \textquotedblleft{}it is allowed\textquotedblright{} (EtymDB; Lewis \& Short). English \emph{licit} and \emph{illicit} come from it.
\end{cousinweb}

\begin{grammarlens}[title={What you've built}]
One more word for telling, explaining and describing at B1.
\end{grammarlens}

\subsection*{Your turn}
\begin{itemize}
\raggedright
  \item \emph{Say it:} \emph{licitus, licita, licitum}
  \item \emph{Say it:} \emph{licitus, licita, licitum}, once more
  \item \emph{Say it:} a sentence of your own with \emph{licitus, licita, licitum}
  \item \emph{Recall:} say the Latin for a rule, then the Latin for lawful, then say \emph{licitus, licita, licitum} again
\end{itemize}

\begin{tcolorbox}[breakable,colback=teal!4,colframe=teal!35!black,title={Before you move on}]
What does licitus, licita, licitum mean? (\textquotedblleft{}Allowed, lawful\textquotedblright{}.) Say it once more.
\end{tcolorbox}
